\documentclass[../main.tex]{subfiles}
\begin{document}
%==========================================TIÊU ĐỀ TẬP========================================
\begin{table}[h]
    \begin{adjustwidth}{-1cm}{}
        \begin{tabular}{>{\centering\arraybackslash}p{6cm}|>{\raggedright\arraybackslash}p{14cm}}
            \multirow{5}{6cm}
            % Cột bên trái: ÉPISODE và số tập
            {\\[-40pt]
                \flaregothic\fontsize{20pt}{20pt}\selectfont \centering
                \color{eptype}HISTOIRE\color{black}\\
                \dawnland\fontsize{72pt}{72pt}\selectfont
                \color{epnum}2
                \\[18pt]
                \flaregothic\fontsize{14pt}{14pt}\selectfont
                \color{epattr}BIENVENUE, \uline{\color{black} Suzuki} !
                % \flaregothic\fontsize{18pt}{18pt}\selectfont
                % \color{epattr}ÉPISODE PERDU
                \color{black}
            }
             &
            % Dòng thứ nhất: tên tập tiếng Nhật
            {\fotskip\fontsize{18pt}{18pt}\selectfont \ruby{虚}{きょ}\ruby{無}{む}に\ruby{飲}{の}み\ruby{込}{こ}まれた！} \\[6pt]
            % Dòng thứ hai: tên tập tiếng Anh
             & {\cafeteria\fontsize{18pt}{18pt}\selectfont The Void}                                       \\[4pt]
            % Dòng thứ ba: tên tập tiếng Việt
             & {\vnmsans\fontsize{18pt}{18pt}\selectfont Khoảng không bí ẩn}                               \\[8pt]
            % Dòng thứ tư: Nhân vật xuất hiện
             & {\caviar{\fontsize{14pt}{14pt}\selectfont Nhân vật: \emp{kuon1} \emp{kaigou1} \emp{suzuki1}}}  \\[8pt]
        \end{tabular}
    \end{adjustwidth}
\end{table}
%===========================================NỘI DUNG==========================================
\color{black}

\pov{Hikawa Kuon}

Tôi không nhớ được mình rơi vào trong đó lúc nào.

Mọi thứ xảy ra quá nhanh, nhanh đến mức tâm trí tôi không kịp phản ứng. Một phần não tôi còn đang cố tính xem nên đến trường muộn bao nhiêu phút, nên giải thích với giáo viên bộ môn như thế nào, phần còn lại đã bị bẻ cong bởi một thứ ánh sáng điên rồ cuốn lấy toàn bộ thân thể.

Chỉ trong chưa đầy một giây, tôi từ giữa lòng thành phố bị hút vào một lỗ xoáy kỳ dị, thứ giống như một miệng cống khổng lồ bằng ánh sáng-- không, nó như được tạo ra từ lỗi phần mềm vũ trụ, chồng lớp ánh tím, ánh trắng, điện khí và bóng tối, tất cả xoắn lại như một tấm vải bị bóp méo.

Và chính khi ấy, tôi sực nhớ ra lời Shion-chan nói, ba ngày trước: ``Nếu đến ngày thứ ba, đừng ra đường. Có thể sẽ có chuyện xảy ra bên ngoài đấy.''

Tôi nhớ đã hỏi lại, lúc đó em ấy còn tỉnh bơ: ``Vài con rối nổ tung thì có gì ghê gớm đâu\dots''

Nhưng bây giờ, khi tôi đang bị cuốn vào giữa một cơn ác mộng thật sự, tôi mới thực sự ngẫm được những lời mà em ấy nói. Những con rối ấy có lẽ chỉ là phần nổi của tảng băng chìm, thứ mà cô phù thủy nhỏ ấy đã vô tình để lại trong thế giới này. Những năng lượng ma thuật cưỡng bức, bản sao không hoàn chỉnh, những thứ vận hành bằng ý chí của người khác\dots\  tất cả là sự can thiệp sai lệch vào trật tự tự nhiên.

Và nếu như vũ trụ có cơ chế tự sửa lỗi, thì cái lỗ xoáy khổng lồ này chính là bản cập nhật lỗi đang xóa sạch những điểm bất thường kia.

\dots Giả thuyết này nghe không khác nào như ở trong các bộ phim khoa học viễn tưởng cả.

Bao nhiêu giả thuyết đều cùng hiện lên cùng lúc, nhưng tôi không thể nghĩ thêm được nữa. Tôi cũng chẳng thể nào biết được đâu mới là thứ đã khiến mọi thứ xoay chuyển nhanh đến vậy chỉ trong chốc lát.

Cơ thể tôi bị xoay tít trong cơn lốc năng lượng. Tôi mất phương hướng. Mất trọng lực. Mất cả nhận thức.

\dots

\dots

\begin{center}
    \uline{The Void.}
\end{center}

*BỤP!*

Cú chạm đất bất ngờ khiến lưng tôi đập mạnh xuống nền.

``Ặc\dots!'' tôi rên lên đau đớn, nhưng ngay lập tức nhận ra một điều kỳ lạ: phía dưới tôi\dots\ không phải mặt đất cứng. Là đệm?

Tôi ngửa đầu nhìn, rồi ngẩng người ngồi dậy.

Đúng thật. Đó là một cái đệm đơn, loại kiểu mà gia đình tôi hay dùng, hơi mỏng, có hoa văn chìm màu be, vải lụa hơi cũ, có mùi\dots\ như của nhà ai đó. Có vết cháy sém nhẹ ở mép dưới, như từng bị xém qua lửa. Cái đệm này có vẻ như\dots\ đã từng nằm trong một căn phòng thực sự. Không, chính xác hơn, nó từng thuộc về một căn hộ nào đó.

Tôi đưa tay ra xoa xoa lưng, vẫn còn cảm giác nhức âm ỉ. Đôi mắt vẫn nhắm nghiền vì cơn đau dội ngược, nhưng tôi cố giữ bình tĩnh. Phải bình tĩnh.

Tôi từ từ ngồi dậy.

Ngay lập tức, tôi cảm thấy rợn người với quang cảnh xung quanh.

Không có gì cả. Không một cái gì hết.

Một màu đen kịt phủ khắp mọi phương. Không có nhà cửa. Không có mặt đất hay bầu trời, chỉ là một không gian vô tận màu đen, mờ mờ ánh tím, giống như đang đứng trong ruột một cái hố không đáy.

Tôi cúi xuống. Bên dưới chân tôi, không rõ là gì, giống như có bề mặt giúp tôi có thể đứng được, nhưng lại không hề phát ra tiếng vang khi tôi vỗ xuống. Cảm giác như đang đứng trên một tấm thảm khí nén.

Trong tay tôi vẫn còn chiếc điện thoại, nhưng máy báo hiệu không có sóng. Nơi này có vẻ như hoàn toàn biệt lập với thế giới bên ngoài.

Tôi cầm lấy cái ba-lô vẫn còn vắt lệch trên vai, kiểm tra lại các giấy tờ và máy móc trong đó. Không mất gì nhiều, cái laptop của tôi vẫn còn nguyên, nhưng tôi biết, nếu có mất\dots\ cũng chẳng ai tới mà tìm lại cho tôi ở cái nơi này.

``\vie{\dots Đây là đâu vậy?}''

Tôi khẽ lẩm bẩm, rồi đứng lên, cố gắng giữ thăng bằng. Đôi chân bước đi chầm chậm về phía trước, dù tôi chẳng biết phía trước là đâu. Mỗi bước đi vang lên trong sự yên tĩnh tuyệt đối, không có gió, không có nhiệt, không có tiếng vang.

Tôi vừa bước đi trong vô định, vừa lẩm bẩm tự hỏi: ``\vie{Mình nên đi đâu bây giờ?}''

Không ai trả lời. Chỉ có tôi, và cái bóng của tôi, ở giữa nơi không thuộc về bất kỳ thế giới nào mà tôi biết.

Một thế giới trống trải, hoàn toàn không có gì cả. Có lẽ tôi sẽ gọi nó là ``\emph{The Void} (khoảng không)''.

\ct

Tôi không rõ mình đã đi được bao xa. Cảm giác về không gian ở đây khác biệt hoàn toàn so với thế giới bên ngoài: không có trọng lực rõ rệt, không có chân trời, và đặc biệt là\dots\ không có điểm đến. Mỗi bước đi như chỉ để giữ cho bản thân không hoảng loạn.

Nơi này cũng không có khái niệm thời gian. Sở dĩ tôi nói vậy, bởi sau khoảng năm phút (ít nhất là theo ước tính trong đầu tôi), đồng hồ đeo tay của tôi vẫn đứng im ở 8 giờ 13 phút sáng.

Tôi giơ màn hình lên để kiểm tra lại, đảm bảo rằng chiếc điện thoại của tôi vẫn ổn.

Màn hình vẫn sáng. Cảm ứng vẫn mượt. Ứng dụng không bị đứng. Nhưng không có sóng. Không có bất kỳ thanh tín hiệu nào xuất hiện ở góc trên.

Tuy nhiên, một điều làm tôi giật mình, là vẫn có tin nhắn mới hiện lên. Người gửi: Yamazu. Có lẽ là tôi mới nhận nó trong lúc tôi đang đuổi theo đồng xu.

Tôi lập tức mở ra đọc: ``\vie{Cô lại bảo vì đau họng nên dời đến tuần sau rồi nhé.}''

Tôi nhìn chằm chằm vào dòng tin nhắn đó, sững người mất vài giây, rồi\dots\ tức điên. Tôi gần như không thể tin nổi mắt mình.

``\vie{\dots Định mệnh, lừa nhau à,}'' tôi nghiến răng, rồi thở phì phò như bò húc.

Tôi chửi không thương tiếc, tuôn ra những lời tục tĩu mà tôi có thể nặn ra được. Thật sự, cái cách mà giáo viên bộ môn của lớp tôi xếp lịch, tôi không biết phải dùng từ gì cho lịch sự nữa.

``\vie{Tuyệt con mịa nó vời lắm! Vừa nãy bảo là đến, bây giờ lại bảo không\dots}'' tôi gầm lên, giọng vang xa rồi lại vọng ngược về từ chính bóng tối vô định.

``\vie{Cô xếp như thế thì bọn em trở tay làm sao kịp!? Lịch báo lúc nửa đêm, rồi sáng ra lại lật kèo, cuối cùng thì em bị hút vào cái thứ của nợ này! \textbf{Toàn tại cô cả đấy!!}}''

Tôi biết mình đang nổi nóng vô ích, nhưng\dots\ tôi phải xả ra. Nếu không thì tôi sẽ phát điên thật sự mất.

Tôi thở ra một hơi dài, rồi khẽ gật đầu, tự trấn an. Chửi mãi thì cũng chẳng giải quyết được gì, mà cũng chẳng có ai bên cạnh để nghe tôi giãi bày, càng chỉ tổ tốn nước bọt.

Tôi mở ứng dụng nghe nhạc trong máy. May mà vẫn hoạt động bình thường. Tôi chọn chế độ phát ngẫu nhiên. Một bài piano nhẹ nhàng vang lên, âm thanh bay ra không có tiếng vọng, như thể được hấp thụ bởi chính hư vô này.

Giai điệu dịu dàng tràn vào đầu tôi, xoa dịu tâm trí đang sôi sục. Tôi cứ để nó phát vậy, nhịp bước của mình khớp theo từng nốt nhạc như một chiếc máy phát nhịp sinh tồn.

Âm nhạc, suy cho cùng, vẫn là liều thuốc tinh thần hiệu quả nhất.

\ct

Thêm khoảng năm phút nữa trôi qua.

Tôi tiếp tục bước, vẫn về một hướng mà tôi cho là ``phía trước''. Nhưng không có gì thay đổi cả, cũng không có một cái chỉ dấu gì để cho tôi biết tôi đang ở nơi nào.

Không cột mốc. Không trạm dừng. Không âm thanh ngoài tiếng nhạc từ tai nghe. Không ánh sáng thay đổi. Thời gian ở đây không trôi, hoặc ít nhất là không vận hành theo cái cách mà tôi từng biết.

``\vie{Và đặc biệt hơn là\dots\ không có người.}'' Tôi bật thốt ra một cách vô thức.

Trong không gian trống rỗng này, giọng tôi vang lên như một làn khói mỏng, rồi tan ngay vào bóng tối. Hẳn nhiên rồi. Làm gì có ai ở cái xó xỉnh này.

Tôi lại cúi đầu, để mặc cơ thể mình đung đưa theo điệu nhạc đang phát.

\dots Nhưng rồi, tôi dừng lại.

Không phải vì bản nhạc dừng. Cũng không phải vì mệt. Mà vì tôi thấy cái gì đó.

Một chiếc xe đạp. Cũ, rỉ sét, bánh trước méo mó, vẫn còn treo biển ``cho thuê 100 yên/giờ.''

Tôi bước đến gần, ngó nghiêng, cố gắng đoán xem\dots\ Nó từ đâu ra?

Không có con đường nào để nó lăn đến đây. Tôi dám chắc là nó cũng đã bị cuốn vào vòng xoáy ban nãy.

Nó không phải là thứ duy nhất tôi thấy. Đi thêm một đoạn nữa, tôi nhìn thấy một biển báo đường kiểu châu Âu, có chữ ``Rue de Strasbourg\footnote{Tên một con đường ở Nantes, một quận giáp biển ở miền Tây nước Pháp.}''. Tôi chưa từng thấy thứ này ở thành phố mình sống.

Những thứ sau đó lại khiến cho tôi nhướng mày, bởi lẽ có những thứ mà tôi không nghĩ chúng có thể tồn tại ở thế giới của tôi.

Một chiếc điện thoại kiểu gập, màu đỏ, nứt vỏ, vẫn sáng màn hình, đồng hồ dừng lại ở thời điểm 3 giờ 24 phút chiều. Nhưng ngôn ngữ hiển thị bằng tiếng Hy Lạp cổ (tôi đoán được vì cấu trúc ngôn ngữ của nó khác biệt so với tiếng Hy Lạp bình thường).

Một con búp bê gỗ Nhật truyền thống, kiểu Kokeshi, nhưng lại được khoác một lớp áo giống như sari của Ấn Độ.

Tôi lặng người.

Cái nơi này không phải là nơi tôi biết. Và những thứ kia\dots\ chúng cũng không thuộc cùng một thế giới. Không hề giống với những vật dụng mà tôi hay bất cứ ai ở thế giới của tôi thường thấy.

Tôi quay vòng quanh, nhìn từng vật thể lơ lửng trong không gian này, đủ hình dáng, từ cờ quốc gia lạ mắt pha trộn giữa cờ Thái Lan với cờ Mông Cổ, hộp bánh in hình nhân vật tôi chưa từng thấy trong đời, đến cả một mảnh báo viết bằng ký tự không phải của bất kỳ hệ chữ nhân loại nào tôi từng đọc qua.

``\vie{Không thể nào\dots}''

Chỉ có một giả thuyết cho tất cả điều này.

``\vie{\textit{Nếu những vật thể này không đến từ cùng một nơi\dots\ thì\dots}}''

Tôi nuốt khan, cảm thấy tim đập nhanh hơn bình thường.

``\vie{\textit{Thuyết đa vũ trụ\dots}}''

Có lẽ\dots\ nó thực sự tồn tại. Thậm chí, tôi còn được tận mắt thấy được nó, dù theo cách gián tiếp, qua những đồ vật mà tôi cả đời chưa từng biết.

``\vie{\textit{Rốt cuộc\dots\ nơi này là gì đây?}}''

Tôi vẫn chưa định hình được giả thuyết về đa vũ trụ trong đầu, giải thích được mọi thứ đang hiện ra trước mắt ra sao thì *BỤP*, tôi va mạnh vào một thứ gì đó.

Không, không phải một thứ. Một ai đó. Chính xác hơn, đầu của một ai đó.

Tôi lùi lại một bước, tay giơ lên bản năng để giữ thăng bằng, điện thoại suýt rơi khỏi tay. Đầu tôi choáng nhẹ vì tưởng đụng phải cục đá.

``\vie{\dots Khoan. Người ư?}''

Tôi chớp mắt liên tục. Là thật. Một người thật sự. Một con người khác ngoài tôi đang hiện diện trong cái thế giới rỗng không này.

``\vie{Đau đau đau\dots\ Cái gì đây\dots}'' cô gái xoa phần sau đầu, rít lên sau cú cụng đầu vô tình kia. Có vẻ như cô ấy cũng nói tiếng mẹ đẻ mà tôi có thể hiểu được.

Người mà tôi vô tình đụng phải\dots\ là một cô gái.

Tôi nhìn chằm chằm. Khoảnh khắc đầu tiên, tôi cứ tưởng đang nhìn thấy hình phản chiếu của chính mình, cho đến khi tôi để ý kỹ hơn.

Cô ấy về cơ bản cao ngang bằng tôi, hầu như những chi tiết trên khuôn mặt cô có nét tương đồng như tôi, nhưng vẫn có chút gì đó để phân biệt được giữa con trai và con gái.

Cô gái ấy mặc một chiếc áo phông gần như y hệt áo tôi đang mặc, chỉ khác màu: tôi đang mặc áo phông xám, còn cô ấy diện áo trắng. Nhưng thay vì quần bò như tôi, cô lại mặc một chiếc váy đỏ đậm, dài tới gần đầu gối, dáng đứng hơi nghiêng sang một bên như thể đang dò xét tôi. Dép quai hậu, giống như tôi đang đi.

Mái tóc đen cô ấy xoăn nhẹ, và được buộc đuôi ngựa phía sau.

``\vie{Giống\dots\ mình à?} tôi thì thầm.

Tôi còn chưa kịp nói gì thì cô gái ấy từ từ quay đầu lại, tay vẫn xoa đầu.

Và tôi ngay lập tức cảm nhận được sát khí. Một thứ gì đó mà Ru-chan hay Sui-chan đều sở hữu, nhưng có phần nguy hiểm hơn nhiều. Thậm chí tôi có thể nhìn thấy cái thứ khí ấy bằng mắt thường.

Đôi mắt xanh nước biển - mang dáng dấp của mẹ tôi - thẫm ánh lên như đèn cảnh báo. Môi cô mím chặt, cổ họng gằn ra như tiếng gầm trầm như con sư tử tìm thấy mồi. Cô ấy xông vào tôi như một con mãnh thú không cần lý do.

``\vie{\textbf{K-Khoan đã! Tôi không--}}''

Một cú đấm như sấm sét bay thẳng vào mặt tôi.

Tôi ngả người tránh sang phải theo phản xạ. Lực gió từ cú đấm rít qua tai, làm tóc tôi tung lên. Nắm đấm của cô đập xuống mặt sàn vô định dưới chân tôi, và gây ra một vết lõm như miệng chảo. Bụi bay khắp nơi, bề mặt thì nứt nẻ tới mức toác ra.

``\vie{\textbf{C-Cái gì thế này!?}}'' tôi sợ hãi lùi lại phía sau. Cô gái ấy không phải người bình thường!

Tôi không biết đây có phải là một ``siêu nhân'' trong các bộ phim hành động phương Tây không, nhưng rõ ràng là có vấn đề với giới hạn sức mạnh ở đây.

Cô ấy không để tôi có thời gian nghĩ hay bất cứ thứ gì để tôi có thể thanh minh. Một tiếng hét vang lên như hổ gầm, tung cú đấm tiếp theo về phía tôi, và tôi lại tiếp tục tránh sang trái.

Cú đấm thứ hai, thứ ba, thứ tư nối tiếp nhau như vũ bão. Tôi phải tránh, lăn, nhảy lùi, và cuối cùng phải dùng cả ba-lô làm tấm chắn tạm thời.

``\vie{Vãi\dots\ Cô ấy mạnh thế\dots}''

Mỗi cú đấm của cô đều nặng như búa tạ, và kỳ lạ thay, với thân hình mảnh khảnh như vậy, điều đó không thể nào hợp lý được. Điều duy nhất mà tô có thể lý giải là do cơ chế thể chất bất thường, hoặc\dots\ bị ``bóp giới hạn'' như trong mấy trò chơi RPG khi chỉnh cheat.

Tôi vừa đổ người, vừa hét lên tuyệt vọng như một con mồi: ``\vie{\textbf{Tôi không biết cô là ai! Tôi không làm gì cả!}}''

Cô gái không hề mảy may quan tâm. Gương mặt đỏ bừng vì giận dữ, miệng cô hét to như thể đang trút cả nỗi đau của một đời người vào tôi: ``\vie{\textbf{Dối trá! Tất cả bọn mày đều là kẻ dối trá!}}''

``\dots Hả?''

``\vie{\textbf{Bọn mày cũng giống nhau cả! Bọn đàn ông chỉ biết cười nhạo trên nỗi đau người khác! Bọn mày thậm chí còn làm tổn thương cả em gái tao nữa, đồ rẻ rách chúng mày!!}}''

Giọng của cô ấy rất uy lực, khá giống với kiểu mà con gái miền Bắc quê tôi hay chửi. Thậm chí là khá sát với phong cách ``giang hồ'' của xứ sở hoa phượng đỏ - \ruby{Kaifan}{Hải Phòng}.

``\vie{Ờm\dots\ Cái gì?}'' tôi vẫn đang ngẩn ra, không hiểu cô ấy đang nói tới cái gì. Tôi thề, cả đời tôi còn chưa gặp cô gái này bao giờ! Tôi cũng chẳng phải như lũ đàn ông thích cười cợt và bức hại phụ nữ - nếu không tính việc tôi ném cờ lê vào mặt của Robo-PON vì phá hủy văn phòng\footnote{Xem lại {\flaregothic ÉPISODE 22}, quyển số Hai.}, nhưng đó là vì cô ấy phá hoại quá mà thôi.

Nhưng câu nói tiếp theo làm tôi điếng người một hồi. ``\vie{\textbf{Chúng mày đã làm gì Suzu rồi!? Khai ra mau!! Em tao đang ở đâu!?}}''

``Suzu?'' Ai cơ? Một cái tên mà tôi còn chưa nghe bao giờ, huống gì là đã từng gặp mặt.

Nhưng tôi không có cơ hội để hỏi lại.

Cô gái bí ẩn kia tiếp tục lao vào. Tôi càng lùi thì cô càng tiến. Cú đánh tiếp theo suýt gạt bay cổ tôi khỏi trục sống. Nếu tôi không dùng phản xạ luyện từ thời tránh đạn giấy của Marine-chan và Mat-chan, chắc giờ tôi đã nằm bất tỉnh trên mặt sàn vô hình kia rồi. Không, có khi tôi còn bỏ mạng ngay tại đây chứ chẳng đùa.

Chẳng mấy chốc, tôi bị dồn đến sát tới đường cùng. Ở đây còn không có tường, không có ranh giới, không có một cái gì để tựa vào, nhưng tôi cảm giác như mình đang bị ép vào góc chết.

Không lối thoát. Mà kể cả có thoát được, tôi cũng không biết chạy đi đâu.

Tôi bắt đầu thở hổn hển. Mồ hôi lạnh rịn ra sau gáy, cảm tưởng cận kề cái chết ngay trước mắt chỉ vì một cái đụng người đầy vô tình.

Cô gái ấy, người đang giơ nắm đấm lên cao kia, sẵn sàng giết tôi bằng tay không, chỉ vì nghĩ tôi là một thằng đàn ông đáng căm thù nào đó, không hề thương tiếc. Một sự hiểu lầm cực kỳ nặng mà tôi còn không có cơ hội để phản biện.

Nhưng ngay khi cô gái định tung cú đấm định mệnh ấy về phía tôi, cô bỗng nhiên khựng lại.

Đôi mắt cô chớp liên tục, như thể đang nghe thấy một giọng gọi từ đâu đó. Ngoài chúng tôi ở đây ra, tôi không nghĩ tôi có thể nghe tiếng nào khác.

``\vie{Chị ở đây mà\dots\ Suzu\dots\ Chị đây\dots}''

Cô gái đứng bật dậy, quay hết bên này đến bên kia, rồi bắt đầu chạy ngược về hướng vô định phía sau. Giọng cô bắt đầu nức nở, vừa gọi vừa khóc: ``\vie{Suzu!? Suzu! Em đang ở đâu? Em đừng bỏ chị mà\dots}''

Tôi đứng đó, ngẩn ngơ, không hiểu chuyện gì đang diễn ra. Vẫn chưa hoàn hồn sau chuỗi đòn vừa rồi, tôi chỉ biết nhìn bóng lưng ấy dần khuất xa, hòa vào khoảng không màu đen ấy.

``\textit{\vie{Không lẽ\dots\ cô ấy đang tìm em gái? Nhưng đây là nơi nào cơ chứ\dots}}''

Tôi nhìn theo dáng hình vừa dữ dội, vừa mong manh ấy. Một cô gái có thể đấm nát mặt sàn bằng thứ sức mạnh vô lý, nhưng cũng có thể gục ngã trước nỗi sợ mất đi một người thân yêu.

Cái thế giới quái quỷ này không chỉ lôi tôi vào để thử thách thể xác, mà có lẽ còn đang đẩy tôi vào một cơn bão cảm xúc chưa từng có trong đời.

Tiếng gọi tuyệt vọng cứ thế vọng dần, vọng dần, rồi sự im lặng bao quanh lấy tôi.

Tới lúc này, tôi mới dám hít một hơi sâu rồi ngồi thụp xuống.

Hai đầu gối vẫn còn run run.

``\vie{Mình\dots\ vẫn còn sống\dots\ Ơn giời.}''

Tôi rút điện thoại ra. Màn hình vẫn sáng. Cảm ứng vẫn hoạt động. Dù không có tín hiệu, mọi thứ trông vẫn ổn, không bị hỏng hóc gì sau khi vừa tránh nạn trong gang tấc.

Tiếp theo, tôi tháo ba-lô, mở ra lấy chiếc laptop mỏng tôi vẫn thường mang theo đi làm và đi học. Mở ra kiểm tra.

Đèn nguồn vẫn sáng. Ổ cứng không kêu tiếng gì kỳ lạ. Quạt vẫn quay đều. Vẫn khởi động bình thường.

Tôi thở phào nhẹ nhõm. ``\vie{May là\dots\ ít nhất đồ vẫn ổn.}''

Không muốn nán lại thêm giây nào nữa ở nơi chiến trường, tôi khoác lại ba-lô, quay đầu chạy theo hướng đối diện mà cô gái kia vừa đi. Tôi không biết mình đang tránh gì, hay định tìm gì, nhưng theo bản năng, tôi muốn tạo khoảng cách với cô gái kia.

\ct

Mười phút trôi qua kể từ trận hỗn chiến một chiều vừa rồi.

Tôi vừa đi, vừa mở lại ứng dụng âm nhạc.

Tai nghe vẫn hoạt động. Một bản nhạc acoustic êm dịu vang lên, lướt nhẹ như sương mai, trái ngược hoàn toàn với cảnh tượng vừa xảy ra chưa đầy ít lâu.

Tôi để mặc mình đi theo giai điệu, như một kẻ du hành lang thang giữa một giấc mơ bị đảo ngược hoàn toàn.

Và rồi, tôi lại bắt gặp những thứ không nên có ở đây.

Một chiếc quạt bàn mini, hiệu Trung Quốc, vẫn còn dây cắm điện lủng lẳng phía sau vào ổ điện. Chiếc quạt vẫn hoạt động, mặc dù tôi dám chắc nơi này không có điện.

Một con chim gỗ sơn kiểu Nga, được khắc hoa văn cầu kỳ, lơ lửng giữa không trung, xoay chậm theo một quỹ đạo vô hình.

Một tấm bản đồ thế giới, nhưng các quốc gia và châu lục lại có hình dạng không giống với Trái Đất tôi biết. Có một lục địa nằm giữa Thái Bình Dương, tên là ``Aetheria''.

Một xe cảnh sát đồ chơi màu xanh biển, nhưng phần logo lại không phải tiếng Anh, cũng chẳng phải tiếng Nhật, mà nó được viết bằng một loại ký tự tròn méo như viết tay trẻ con. Làm tôi nhớ đến vụ cải cách giáo dục mà bọn trẻ con đọc mỗi ô vuông hay hình tròn là một từ, sau này tôi mới biết đó là cách đếm tiếng.

Tôi dừng bước, nhìn từng thứ một. ``\vie{\textit{Không phải là một thế giới. Là nhiều thế giới}},'' tôi thì thầm, như đang thốt ra kết luận sau một chuỗi thí nghiệm.

Mỗi vật thể ở đây là một mảnh nhỏ của các chiều không gian khác nhau, và tôi, không hiểu vì lý do quái gì, lại rơi vào đúng cái mớ hổ lốn đó, chỉ biết rằng tôi đã vô tình trở thành một phần trong số chúng, xét theo khía cạnh nào đó.

Dù vậy\dots\ vẫn có điều gì đó níu tôi lại với thực tại.

Tôi vẫn không ngừng nghĩ về cô gái lúc nãy.

Cô ấy là ai?

Tại sao cô ấy lại mạnh đến vậy?

Và tại sao cô ấy lại phản ứng dữ dội đến thế khi nhìn thấy tôi?

Có lẽ là vì một ai đó vô tình thốt ra chăng? ``Suzu'', đúng không nhỉ?

Có lẽ là em gái của cô ấy chăng?

Nghe giọng nói lạc đi trong nước mắt ấy\dots\ tôi không nghĩ cô đang diễn. Lại càng vô lý hơn khi cô ấy diễn ở một nơi mà tôi không rõ nó nằm ở đâu, ở bất kỳ chỗ nào mà tôi từng biết.

Cô thực sự hoảng loạn. Thật sự rất hoảng.

Một người có thể tung cú đấm làm sụp mặt đất, nhưng cũng có thể khóc như một đứa trẻ khi sợ mất người thân.

Sự mâu thuẫn ấy\dots\ làm tôi nhớ đến ai đó. Như Subaru-chan, hay thậm chí là\dots\ chính tôi. Nếu tôi ở trong tình cảnh đó, có lẽ tôi cũng sẽ khản cổ gọi người thân như vậy.

Tôi vẫn không biết cô ấy tên gì. Cũng chẳng rõ tôi sẽ gặp lại cô ấy hay không.

Nhưng cảm giác ấy, cái cảm giác chạm trán với một người giống tôi đến kỳ lạ, cả về ngoại hình lẫn cơn giận dữ âm ỉ trong lòng, cứ lặp lại mãi trong đầu tôi, cùng với tiếng nhạc đang khe khẽ vang lên bên tai.

Dù sao thì tôi phải tìm cách rời khỏi nơi ma xui quỷ hờn này đã.

\ct

Tôi không biết mình đã đi thêm bao lâu. Có lẽ là năm phút nữa. Có thể là mười, hoặc mười lăm, hai mươi\dots\ Nhưng trong cái nơi gọi là ``khoảng không'' này, thời gian gần như vô nghĩa.

Mỗi bước chân lướt qua một vùng tối không màu, không âm vang. Vẫn chỉ là tôi và âm nhạc, và hàng loạt những mảnh vụn vũ trụ cứ lần lượt xuất hiện trước mắt.

Ban đầu là những thứ nhỏ nhặt, những thứ mà tôi vừa thấy. Nhưng càng đi, tôi càng gặp những thứ\dots\ lớn hơn. Kỳ lạ hơn. Không thể nào giải thích nổi.

Tôi thấy một cửa hàng tiện lợi, đúng nghĩa là một cái cửa hàng đầy đủ, chỉ thiếu phần nền móng. Nó lơ lửng cách mặt ``đất'' chừng ba mét, toàn bộ đồ đạc bên trong vẫn còn nguyên: giá kẹo, máy tính tiền, cửa kính vỡ một nửa. Không có người.

Tiếp đó là một toa tàu điện ngầm, bị cắt làm đôi không biết bằng cách nào, nằm xoay nghiêng trong không trung. Bên trong vẫn còn một chiếc ghế có vết ngồi lõm, một túi xách lạc lõng nằm trên sàn kim loại đã gỉ.

Đi tiếp một đoạn nữa, tôi bắt gặp một khối cầu khổng lồ có kết cấu như một con mắt, xoay nhẹ như đang quan sát tôi. Không rõ là sinh vật hay máy móc. Tôi nín thở mà bước qua, không dám hỏi vì sợ nó nghe thấy.

Tôi đã từng trải qua những thứ kỳ quặc và quái đản ở văn phòng \hll, nhưng đến mức thế này thì\dots\ quả đúng là làm tôi nổi da gà.

Tôi bắt đầu nhận ra: những gì xuất hiện ở đây không chỉ là vật thể vô tri. Một vài thứ có thể là tàn dư sống sót từ các thế giới từng sụp đổ.

Giữa không gian kỳ quặc ấy, tôi nhìn thấy một chiếc bàn ăn kiểu Nhật.

Một cái bàn không lớn, chỉ là loại bàn tròn thấp, dùng cho ba bốn người, kiểu gia đình nhỏ, ước chừng gia đình tôi. Trên mặt bàn vẫn còn vương vãi những chiếc đĩa sứ nhỏ, một vài chén nhựa màu hồng và trắng, và một cái nồi nhôm con con vẫn còn mở nắp.

Tôi tò mò tiến lại gần. Cảnh tượng trước mắt giống như một bữa ăn bị bỏ dở, hoặc đúng hơn, là bữa ăn đã kết thúc cách đây không lâu.

Các đĩa đã sạch trơn. Hầu như không còn lấy một thứ gì. Chỉ có vài hạt cơm dính lại ở mép bát, một cái thìa cong nằm nghiêng, và một vài mảnh bánh tempura bị cắn dở, bị gặm đến sát vỏ như thể bị ``xử lý'' cực kỳ tận tình. Chiếc nồi canh cũng vậy, hoàn toàn không còn lấy một giọt nước nào. Chỉ còn vài vệt dầu, và một mùi thơm nhè nhẹ của nước hầm miso. Tôi cúi đầu xuống ngửi, mùi umami của nấm xộc thẳng vào mũi. Khá ngon lành.

Một ai đó đã ăn rất nhanh, rất kỹ, như đã trải qua cả ngày dài không có gì vào bụng. Không thể là một con thú, vì chằng có động vật nào lại có thể ăn sạch bách như thế.

Tôi khẽ nhếch môi: ``\vie{Ai mà ăn như lợn thế này\dots}''

Nhưng câu nói vừa bật ra đã khiến tôi chững lại.

Một cái bàn ăn giữa hư không.

Một người từng ngồi đây, ăn tất cả mọi thứ, không để lại gì, trừ vài hạt cơm dính và vệt mỡ mỏng.

``\vie{\textit{Thật là vãi chưởng\dots}}''

Tôi bắt đầu cảm thấy có ai đó khác đang ở gần. Cô gái với sức mạnh như siêu nhân kia là người duy nhất tôi gặp.

Nhưng tôi cảm tưởng một người khác đã làm điều này. Có lẽ\dots\ người đó đang ở rất gần.

\ct

Tôi lang thang thêm một đoạn khá xa nữa, vẫn là bước đi không có điểm dừng, tay đút túi quần, tai nghe phát ra những giai điệu quen thuộc từ bài nhạc cuối cùng trong danh sách. Một hòa âm nhẹ nhàng mà tôi thường nghe khi ngồi một mình ở ban công, giờ đây lại vang lên giữa một nơi mà tôi không biết là liệu có đúng với tình cảnh hiện tại của tôi hay không.

Tôi lẩm bẩm theo giai điệu, câu chữ rơi ra vô thức, giống như một bản phản xạ tự nhiên. Ở một nơi mà không có lấy nổi một bóng người, tôi tự hỏi mình còn giữ bao nhiêu phần bình thường trong cái bất thường này. Không biết bao lâu nữa tôi sẽ phát điên ở một nơi mà cả không gian lẫn thời gian đều hoàn toàn đảo lộn.

Vừa dứt lời cuối cùng, bài hát cũng tắt. Danh sách phát kết thúc.

Tôi lấy điện thoại ra, vuốt sang một album khác. Một danh sách nhạc từ trò chơi điều khiển bóng của Tàu Khựa mà tôi hay chơi khi rảnh. %hint: Rolling Sky

Nhưng ngay lúc tôi định nhấn nút phát thì\dots

``\vie{\dots Onee\dots\ sama\dots}''

Một tiếng gì đó rất khẽ, thoảng như gió. Giọng nói này không giống với cô gái ban nãy.

Tôi khựng lại.

Tai tôi nghe rõ. Một âm thanh\dots\ của tiếng con gái. Một cô bé đang khóc.

``\vie{\textbf{\dots Onee\dots\ sama\dots!!}}''

Tiếng đó lại vang lên, lần này có phần nức nở, thổn thức, như thể nghẹn trong cổ họng. Không lớn, nhưng có sức nặng lạ kỳ.

Tôi tháo tai nghe, lắng nghe lần nữa. Tiếng khóc ấy vọng từ hướng Tây, quãng không xa lắm.

``\vie{Chị đang ở đâu\dots\ Cứu em với\dots}''

Tôi định sẽ phớt lờ. Gặp một người con gái ở nơi này lần trước suýt nữa khiến tôi mất mạng. Tôi không muốn chuốc họa vào thân một lần nào nữa.

Tôi chẳng phải là một anh hùng luôn cứu mỹ nhân. Chẳng có lý do gì để tôi có thể tự đưa mình vào thế nguy hiểm một lần nữa, chỉ để nhận lấy cái chết bất ngờ cả.

Nhưng rồi\dots

``\vie{\dots Onee-sama\dots\ em sợ lắm\dots\ Mau tới cứu em đi\dots}''

Tiếng gọi ấy cứ thấm vào người tôi, như có gì đó không đúng. Nghe như một đứa trẻ đang lạc mất người thân trong cơn hoảng loạn.

Có vẻ như cô gái ban nãy không phải là người duy nhất thất lạc ở nơi này. Phải chăng người đó với cô gái ban nãy là cùng một gia đình? Cũng có thể lắm chứ.

Tôi thở dài, nhìn về phía Tây, rồi bước đi, lặng lẽ nhét tai nghe vào túi áo.

\ct

Cảnh vật dần đổi khác khi tôi rẽ vào một đoạn không gian sẫm màu hơn. Những khối hộp trôi lơ lửng bên trên, một vài bụi cây nhỏ không rõ nguồn gốc run rẩy, mặc dù tôi không nghĩ nơi này có gió. Tôi cảm thấy lạnh sống lưng, nhưng không rõ tại sao.

Bước thêm vài bước nữa, tôi đã nhìn thấy người đó.

Một cô bé có dáng người nhỏ nhắn, ngồi thu mình lại gần mép một vách nứt cạn, hai tay vòng lấy đầu gối, mái tóc đen dài tới lưng, được cột hai đuôi ngựa thấp. Trên đỉnh đầu là một chỏm tóc ngố bật lên, hơi ngả sang trái, khẽ chõng xuống như đang rầu rĩ.

Từ xa nhìn lại, cô ấy giống hệt em gái tôi, Ryuuhou. Nếu không vì chỏm tóc ngố đó, có lẽ chắc tôi đã nhận nhầm em ấy mất.

Tôi chững lại.

``\vie{Không\dots\ Không thể nào\dots}''

Tôi nhìn kỹ hơn. Cô bé mặc một chiếc áo dài tay màu lá mạ, bên dưới là chiếc đầm đỏ sẫm dài qua đầu gối, đi cùng quần tất đen và giày vải vàng đơn giản. Có gì đó khác lạ.

Đôi mắt của cô bé lại là một thứ gì đó ấn tượng.

Một bên xanh lam, có cùng màu với đôi mắt của cô gái hồi nãy. Một bên đỏ thẫm giống như tôi. Giống hệt em tôi, nhưng ngược vị trí hoàn toàn. Ryuuhou có mắt trái xanh, mắt phải đỏ. Còn cô bé này thì ngược lại.

Tôi khẽ rùng mình. Không phải Ryuuhou. Nhưng là ai đó rất giống.

Tôi bước lại gần, cố nhón chân nhẹ nhất có thể.

``\vie{\hl{N-Này\dots}}'' tôi mở lời, khẽ hạ giọng. ``\vie{\hl{Em ổn chứ\dots?}}''

Cô bé giật mình, ngẩng đầu lên nhìn tôi, và ngay lập tức lùi lại một bước, ánh mắt hoảng loạn. Cô bé định mở miệng ra hét, nhưng dường như không phát được ra tiếng, có lẽ vì quá sợ.

Tôi khựng lại ngay, hai tay giơ lên ngang tầm vai để thể hiện mình không có ác ý. \vie{Không sao, anh không làm gì em đâu\dots\ đừng sợ\dots}

Cô bé vẫn giương đôi mắt sợ sệt nhìn tôi đầy nghi ngờ, mắt không rời khỏi tay tôi lấy một giây. Lưng tựa sát vào mặt đất như để thủ thế bỏ chạy.

Cũng dễ hiểu thôi. Một người bất ngờ xuất hiện giữa một nơi khỉ ho cò gáy thế này, việc bị giật mình cũng không lấy làm lạ.

Nhưng khi tôi nhìn vào đôi mắt long lanh đang chực chờ nước mắt kia, cô bé đang cho tôi một cảm giác khác. Không phải là nỗi sợ vì gặp người lạ như thông thường.

Mà là một thứ gì đó khiến cho cô bé ám ảnh hơn nhiều.

Tôi nuốt khan. ``\textit{\vie{Cô bé này\dots\ sợ đàn ông đến vậy sao?}}''

Tôi đoán có thể là do bị bắt nạt, hoặc do một người anh, hoặc một người chị quá bảo vệ? Có thể cô ấy là ``Suzu'' mà cô gái hung hăng trước đó đã nhắc đến?

Tôi ước rằng tôi có thể hỏi thẳng cô bé kia, nhưng với tình cảnh hiện tại, hầu như chỉ tiếp cận gần thôi là rất khó.

Tôi còn đang nghĩ xem tôi đang định làm gì để tiếp cận cô bé ấy, thì bỗng nhiên\dots

*ỌC ỌC*

Tiếng bụng réo lên.

Cả hai chúng tôi đều im bặt. Tôi nhìn hướng theo cô bé. Cô bé đỏ mặt, cúi gằm xuống vì xấu hổ.

``Hầy\dots''

Tôi thở dài. Tình huống này không lạ với tôi. Em gái tôi cũng từng như thế. Bình thường, nếu là Ryuuhou, tôi có thể cho tiền, hoặc để tự em ấy xài tiền mua đồ ăn. Nhưng ở một nơi mà hầu như không có gì ngoài những vật dụng kỳ quặc bay lơ lửng khắp nơi, việc đó là bất khả thi.

Tôi thoáng nhớ lại cái cửa hàng tiện lợi mà tôi đã đi qua trước khi bắt gặp cô gái hung hăng ban nãy, nhưng rồi khẽ lắc đầu. Nơi này còn không có biển chỉ dẫn, làm sao tôi biết được là nên đi đâu? Chưa kể nếu tìm lại được cửa hàng đó, cũng không chắc là nó sẽ có cái gì để mua đâu.

``\vie{Em đói đúng không?}'' tôi cất giọng hỏi. Cô bé ấy gật đầu, hai tay vẫn ôm lấy bụng.

``\vie{Được rồi, chờ anh chút.}''

Tôi cẩn thận lôi ra hộp cơm trưa từ ba-lô. Một hộp cơm nhà làm, đơn giản thôi: cơm nắm cá ngừ, trứng cuộn, vài lát xúc xích và rau luộc. Không có gì đáng ngờ. Mọi khi tôi sẽ đi ăn ngoài hoặc sử dụng lại đồ thừa từ hôm qua, nhưng hôm nay tôi muốn làm một điều gì đó khang khác so với bản thân, và thế là tôi đã tự tay làm hộp cơm này từ sáng sớm, với nguyên liệu tươi mua ở chợ.

Tôi đặt hộp cơm xuống, rồi ngồi lùi lại.

``\vie{Anh sẽ để ở đây. Tin anh đi, anh không bỏ gì kỳ lạ vào cả. Nhà anh làm đấy,}'' tôi gượng cười, không biết liệu cô bé có tin lời tôi nói không.

Để cho chắc chắn hơn, tôi đã phải buông một lời nói dối - mà thực chất là nói thật: ``\vie{Anh cũng ăn rồi. Không có thuốc độc đâu.}''

Cô bé nhìn tôi. Rồi nhìn hộp cơm. Rồi nhìn tôi lần nữa, nét mặt vẫn hiện rõ vẻ do dự. Tôi không ngạc nhiên.

Sạch là một chuyện, nhưng vẫn còn một vấn đề khác mà tôi lo hơn: liệu hộp cơm của tôi có khẩu hợp khẩu vị với cô ấy không?

Để rồi, nỗi lo ấy của tôi xem chừng là hơi thừa.

Sau khoảng ba giây, cô bé nhào tới, gắp lấy đũa và thìa từ hộp cơm mang theo, rồi ăn ngấu nghiến tới mức tôi toát mồ hôi hột.

Cơm nắm tan biến trong tích tắc. Trứng cuộn chỉ kịp nghe tiếng ``chẹp'' là mất hút. Miếng xúc xích không kịp nhìn kỹ đã không còn. Mọi thứ trong hộp cơm bị xử lý sạch sẽ trong vòng chưa đầy một phút.

Tôi há hốc mồm. Chưa bao giờ thấy ai ăn nhanh như vậy. Như thể cô bé đã bị bỏ đói mấy ngày nay.

``\vie{C-Chờ đã\dots\ Ă-ăn chậm thôi! Không ai tranh của em đâu mà sợ!}'' tôi vội lên tiếng, nhìn thấy cô bé vừa ăn vừa nuốt với tốc độ đáng kinh ngạc.

``Khụ khụ\dots\ Ặc\dots'' cô bé mắc nghẹn, ho từng hồi dài.

``\vie{Biết ngay mà\dots\ Đây!}'' tôi vội vàng lôi chai nước ra đưa. Một chai nước tôi lấy từ bình lọc ở nhà, hoàn toàn sạch.

Cô bé liền cầm lấy ngay, tu một hơi dài đến cạn.

Tôi vẫn chưa hết bàng hoàng. Tôi không nghĩ cô bé đói lả tới mức thế, mà đúng hơn\dots

``\vie{\textit{\dots chẳng phải như thế này là quá phàm ăn rồi đó sao?}}''

Nhưng trong lúc tôi còn đang suy luận xem tại sao cô bé bí ẩn kia lại ăn nhanh như vậy, tôi ngó xuống hộp cơm mà chỉ mới ú ụ một phút trước đó, giờ đã hết nhẵn không còn lấy một thứ gì. Chai nước mà tôi mang theo cũng không còn.

``\vie{\hl{Cảm\dots{} ơn\dots{} anh,}}'' cô bé thỏ thẻ, mắt nhìn xuống đất, giọng nhỏ như tiếng muỗi kêu. Chất giọng của cô có hơi hướng gần giống như Ryuuhou, nhưng cao hơn một chút.

Tôi khẽ mỉm cười, thở ra nhẹ nhõm. Có vẻ như\dots\ cuối cùng cũng tìm thấy một người không muốn giết tôi ở nơi này.

\ct

Cô bé ấy ngồi lại trên nền trơn lạnh của không gian tối mịt ấy một lúc lâu. Sau khi uống cạn chai nước, cô vẫn không nói gì thêm, chỉ khẽ cúi đầu, giữ ánh mắt tránh xa tôi, dù vẻ mặt đã bớt hoảng loạn hơn trước.

Tôi ngồi đối diện, tay xoa nhẹ gáy, vừa để thư giãn, vừa quan sát cử chỉ cô bé.

``\vie{\dots Em tên là Suzuki-chan, đúng không?}'' tôi hỏi, giọng hạ thấp nhất có thể. Tại sao tôi lại biết được tên em ấy, chắc bạn đang tự hỏi?

Đó là vì ngay sau khi cô bé cảm ơn tôi về hộp cơm trưa, em ấy đã tự nói tên của mình ra. Suzuki (\ruby{鈴}{すず}\ruby{木}{き}, \ruby{Xư-dư}{Linh}-\ruby{ki}{Mộc}), với tên trộn lẫn từ họ của Nene-chan và cái cây.

Cũng có thể tên này có quan hệ nào đó với cái tên ``Suzu'' mà cô gái trước đó có nhắc tới, từ đó tôi đang nghiêng về giả thuyết là cô bé này với cô gái trước là hai chị em. Dù vậy, tôi vẫn không hoàn toàn chắc chắn với điều đó.

``\vie{Anh có nghe thấy em gọi `Onee-sama' nãy giờ. Không lẽ\dots\ em bị lạc chị mình?}''

Suzuki-chan khẽ gật đầu. Một cái gật, rất nhỏ.

Tôi khẽ thở ra, lòng nhẹ bẫng. Dù chỉ là một động tác bé nhỏ, nhưng với một cô bé nhút nhát và cảnh giác cao độ như vậy, đó là cả một bước dài.

``\vie{Anh là Kuon,}'' tôi mỉm cười, tay chỉ vào ngực mình, như thể đang chỉ dẫn một đứa trẻ sơ sinh mới lọt lòng bằng giọng nhẹ nhàng nhất có thể. ``\vie{Em không cần phải sợ đâu. Anh không phải kẻ xấu.}''

Cô bé không đáp, chỉ cúi đầu thấp hơn.

Tôi chẳng trách được. Nếu là em gái tôi bị kẹt ở nơi quái đản thế này mà gặp người lạ, có lẽ con bé cũng sẽ phản ứng như thế. Nghĩ đến đó, tôi lại bất giác nhìn cô bé thêm lần nữa.

Suzuki-chan trông giống Ryuuhou thật. Không hoàn toàn giống, nhưng từ ánh mắt, khuôn mặt nhỏ, và cả cái cách tay ôm lấy gối khi ngồi thu mình lại, đều khá tương đồng. Và tôi không thể làm ngơ được.

``\textit{Nếu mình không bảo vệ được ai ở đây, thì ít nhất\dots\ cũng phải bảo vệ được cô bé này.}''

Tôi ngồi thẳng dậy, gật nhẹ đầu một cái như để xác nhận lại với bản thân.

``\vie{Anh nghĩ nên tìm Onee-sama của em cùng nhau,}'' tôi ngập ngừng, rồi hỏi tiếp, ``\vie{vậy\dots\ em muốn đi cùng anh không?}''

Suzuki-chan ngẩng đầu nhìn tôi lần đầu tiên từ lúc cả hai gặp nhau. Đôi mắt xanh-đỏ ấy ánh lên một tia lưỡng lự.

Tôi đoán được điều cô sắp nói. Và tôi lên tiếng trước.

``\vie{Anh biết. Chắc chị em đã dặn em `không được đi theo người lạ' rồi, đúng không?}''

Cô bé giật mình, khẽ gật đầu.

``\vie{Anh không trách em. Nếu là Ryuuhou - em gái anh - thì anh cũng sẽ dặn nó như thế,}'' tôi cười nhẹ, ``\vie{nhưng nơi này đâu còn ai khác, phải không?}''

Suzuki-chan nhìn quanh. Không một bóng người. Không một âm thanh. Chỉ có một khoảng không vắng lạnh.

``\vie{\hl{\dots Vậy\dots{} em sẽ đi cùng anh\dots{} dù chỉ một chút thôi ạ,}}'' cô bé nói nhỏ, giọng lí nhí như tiếng gió.

Tôi mỉm cười, vươn tay về phía cô bé.

Cô ngập ngừng một chút, rồi khẽ đưa tay ra nắm lấy tay tôi.

Cảm giác tay cô bé nhỏ và lạnh, run nhẹ, nhưng cũng rất chắc. Tôi siết tay lại một cách vừa phải, rồi nhẹ nhàng đứng dậy, kéo cô bé cùng bước.

\ct

Trên con đường ấy, khi chúng tôi tiếp tục đi về phía vô định kia, tôi có liếc sang và hỏi, ``\vie{Em có muốn anh cõng không? Đi bộ thế này chắc mỏi chân lắm.}''

Suzuki-chan thoáng giật mình. Nhưng rồi lắc đầu. ``\vie{\dots Em ổn ạ. Em vẫn đi được.}''

Tôi không gặng thêm. Tôi tôn trọng điều đó. Nhưng đâu đó trong ánh mắt cô bé, tôi thấy có thứ gì đó phức tạp hơn. Một điều gì đó\dots\ khó nói hơn thế.

\pov{-}

Cô bé tóc hai bím vẫn đang líu ríu đi theo cậu con trai với những bước nhỏ và nhẹ như bông. Một tay nắm lấy tay của Kuon, một tay khẽ nắm lấy phần dưới của mình.

Suzuki thật sự đã rất mỏi chân. Nhưng cô không muốn nói ra.

Không phải vì ngại, dù đúng là cô vẫn còn e sợ Kuon.

Không phải vì tự ái, cô chưa từng là người giỏi chịu đựng.

Mà là vì\dots\ cô bé không muốn làm phiền.

Ba ngày qua, cô đã lang thang một mình ở nơi quái lạ này. Nơi mọi thứ dường như bị cắt đứt khỏi thế giới thật. Thức ăn thì hiếm. Nước còn hiếm hơn. Và cô cũng chẳng dám đi vệ sinh.

Từ hồi còn sống trong khu chung cư nghèo, cô đã quen với việc ``nhịn''. Khi chỉ có một chai nước cho cả ngày, Suzuki luôn cố gắng uống ít nhất có thể. Rồi từ đó, cũng quen luôn với việc nhịn luôn cả những chuyện mà người ta vẫn cho là bình thường. Từ việc ăn uống, tắm giặt, cho tới vệ sinh cá nhân, tất cả đều chỉ dừng lại ở mức tối thiểu nhất, mặc cho Kaigou cứ giục mãi, nhưng đáp lại chỉ là một cái cười nhỏ nhẹ từ cô bé: ``Không sao đâu, chị ạ. Em vẫn có thể chịu đựng được.''

Để rồi giờ đây, Suzuki đã lỡ uống quá nhiều nước, cả canh lẫn chai nước Kuon đưa, đến mức bụng của cô căng lên vì chứa nhiều nước, cộng với việc cô chưa được đi vệ sinh từ khoảng thời gian dài đó, đã khiến cho cô có phần hơi ngượng khi đề cập tới nó.

Nhưng rồi\dots\ trước người con trai mà cô chưa từng gặp kia, người mà đã giúp cô vô điều kiện mà không có động cơ thâm độc gì, thậm chí còn thẳng tính thừa nhận, rồi còn cùng rủ cô đi tìm chị gái thất lạc của mình mà không cần trả ơn, cô bé lại cảm thấy cảm kích.

Cô nở một nụ cười nhỏ, ánh mắt ngước lên cậu con trai còn đang quay quanh tìm đường.

``\textit{\vie{Được đi bên cạnh anh\dots\ em cảm thấy an tâm hơn hẳn\dots\ Kuon-nii-san ạ.}}''

\pov{Hikawa Kuon}

Tôi không nói gì thêm, cũng không hỏi thêm.

Chỉ biết, ở nơi mà tôi không có bản đồ, không có phương hướng, cũng chẳng có gì chắc chắn, thì việc có ai đó đi cùng, dẫu chỉ là một cô bé ước chừng kém tôi mười tuổi, cũng đã là một điểm tựa quý giá.

Tôi nắm nhẹ tay cô bé hơn, rồi nhìn lên con đường phía trước: một màu đen lặng thinh, không trần, không nền, không điểm đến.

``\vie{Đi thôi. Chúng ta sẽ tìm được lối ra, hoặc ít nhất\dots\ là tìm lại được chị của em đã.}''

``\vie{Vâng.}'' Một câu trả lời dù vẫn còn khá nhỏ, nhưng có một chút tự tin hơn, và đó là điều cảm thấy an tâm hơn.

\ct

\dots

\dots

Tôi không biết mình và cô bé ấy đã đi được bao lâu, chỉ biết rằng mỗi bước chân đều là những bước mơ hồ giữa vùng đất không phương hướng. Âm nhạc trong tai tôi vẫn tiếp tục vang lên, nhưng lần này, tôi để một bên tai nghe trượt xuống, để tiện nghe xem Suzuki-chan có nói gì không.

Có vẻ sau khi ăn no và uống đủ nước, cô ấy bắt đầu cởi mở hơn. Cô bước cạnh tôi, dáng đi vẫn có phần rụt rè, nhưng đã không còn nhìn tôi với ánh mắt phòng thủ như trước.

``\vie{Kuon-nii-san này\dots}'' giọng cô bé nhỏ nhẹ cất lên, gần như thì thầm, ``\vie{anh có thấy một\dots\ cái điện thoại không? Màu đen, viền trầy trầy, kiểu dáng cũ\dots\ đời 2018 ấy ạ\dots}''

Tôi hướng về phía mà cô bé chỉ, thứ mà tôi đang cầm trên tay. Một chiếc điện thoại cùng hãng với máy tôi đang dùng, nhưng đời trước đó khoảng hai năm. Một máy dòng A đời 2018.

Các bạn đang tự hỏi tại sao tôi lại có chiếc máy này? Đó chính là chiếc điện thoại mà cô gái kia vô tình đánh rơi khi đang hành hạ tôi ban nãy. Tôi nhặt máy đó lên và định sẽ tìm cách để trả lại cho cô ấy, nhưng với tình hình hiện tại, tôi còn chẳng biết liệu tôi có thể gặp lại cô ấy được không nữa.

``\vie{Anh biết loại này. Sao vậy?}''

Cô bé khẽ gật gù. ``\vie{Em nghĩ là\dots\ Onee-sama đánh rơi nó đâu đó khi bị hút vào đây\dots}''

Tôi hơi nhướng mày, vội cúi xuống nhìn chiếc điện thoại đang nằm trong lòng bàn tay mình. Một thoáng giật mình thoáng qua, nhưng tôi kịp nuốt lại, không dám để lộ. Tôi không nên vội vàng kết luận mình từng đối mặt với một người có thể là Onee-sama của em ấy, ít nhất là chưa phải lúc này.

Tôi vô thức hỏi ngược lại cô bé: ``Onee-sama?''

``\vie{Vâng\dots}'' cô bé gật nhẹ. ``\vie{Bọn em đến từ một nơi khác. Trước khi bị hút vào đây, bọn em sống trong một chung cư. Căn hộ nhỏ thôi, nhưng ấm áp\dots\ vì có chị ở bên.}''

Nghe vậy, tôi lặng người.

Căn hộ\dots\ Chưng cư\dots\ Có thứ gì đó khiến tôi nhớ đến tấm đệm mà mình đáp xuống khi rơi vào nơi này. Hỏi em ấy một câu về nó\dots\ chắc là không chết ai đâu nhỉ?

``\vie{Mà này\dots}'' tôi quay sang, ``\vie{em có biết một cái đệm mỏng, có hoa văn chìm màu be, bị cháy sém ở góc không? Lúc anh rơi xuống đây, anh đập ngay vào cái đệm đó.}''

Cô bé chớp mắt. ``\vie{\dots Đệm có hoa văn chim màu be?}''

Tôi gật.

``\vie{\dots Cái đó\dots\ là của bọn em ạ.}''

``\vie{Thế ư?}''

``\vie{Vâng\dots\ Em và Onee-sama vẫn nằm đệm đó,}'' cô nói, rồi chợt ngập ngừng. ``\vie{À, còn cái vệt cháy ấy\dots}''

Tôi nuốt nước bọt, chờ cô bé tiếp tục. Cô cúi đầu thấp xuống, giọng nhỏ hơn hẳn.

``\vie{Chung cư bọn em\dots\ từng bị cháy. Lúc ấy em sợ lắm, khói mù mịt\dots} Suzuki bủn rủn tay chân, rồi tiếp tục: ``\vie{Nhưng chị đã kéo em chạy trước, không mang theo được gì cả. Sau đó, em không rõ tại sao một số đồ lại xuất hiện ở đây\dots}''

Tôi thở một hơi dài. Quả đúng là một trải nghiệm kinh hoàng mà cô bé đã phải trải qua, và vì thế, rõ ràng không dễ để kể lại, nhất là với một cô bé vừa mới bắt đầu tin tưởng người lạ. Tôi im lặng một lúc, để cơn căng thẳng trong cô bé tự tan đi.

Tôi quyết định chuyển chủ đề khác bớt xám xịt hơn. ``\vie{Lúc anh lảng vảng ở đây, anh có thấy một bàn ăn bốn người, có mấy cái chén bát vương cơm vãi ra. Em biết cái đó không?}''

Cô bé ngẩng lên, mắt sáng hẳn. ``\vie{A, biết ạ! Là do em ăn đấy ạ!}''

``\vie{\dots Em ăn á?}'' tôi thoáng ngạc nhiên.

``Vâng!'' Suzuki-chan gật đầu lia lịa. ``\vie{Lúc đó em đói, thấy có đồ ăn thì không kìm được\dots\ Nhưng mà, bàn đó nhiều thật luôn! Em ăn hết chắc cũng phải ba phút đấy! Súp miso nấm ấm bụng lắm, rồi bánh tempura giòn rồm rộp nữa chứ!}''

Giờ tôi hiểu tại sao cái bàn ăn đó lại không còn gì rồi.

Tôi nửa như muốn phì cười, nửa ngỡ ngàng với lời bộc bạch của cô bé. ``\vie{Ba phút mà dọn sạch được cái bàn ấy á?}''

Cô bé khẽ phồng má, như thể bị trêu: ``\vie{Em ăn nhanh lắm đó! Nhưng mà, từ bé em đã thế rồi, ăn bao nhiêu cũng không béo.}''

Tôi hơi nghiêng đầu, không tin nổi vào tai mình: ``\vie{Không béo thật á? Anh thấy em ăn kiểu đó thì ít người tin lắm đấy.}''

``\vie{Thật mà! Hồi xưa em còn ăn nhiều hơn mà vẫn gầy queo à. Onee-sama cứ bảo em là `cái động không đáy biết đi',}'' Suzuki hồ hởi, cười khúc khích với câu chuyện của bản thân.

Nụ cười đó không hề gượng gạo, mà đó là một cảm xúc thật. Một nụ cười hồn nhiên của một cô bé vừa bước sang tuổi thiếu niên, hiện rõ sự trong sáng trong cái không gian u ám tĩnh mịch này.

Tôi cũng bật cười theo, vừa thấy câu chuyện có phần hài hước, khi cô bé này giống Ryuuhou ở vài điểm: cũng nhỏ nhắn, cũng dễ thương, và cũng ăn mãi không thấy lên cân, nhưng vừa cảm nhận được sự chân thành của mình.

Dẫu vậy, trong từng lời nói và ánh mắt của cô bé vẫn còn một lớp dè chừng, một khoảng cách mà tôi chưa chạm tới được. Nhưng ít ra, tôi nghĩ\dots\ tôi đã tiến được một bước xa hơn rồi.

Chúng tôi cứ thế tiếp tục đi. Tôi mỉm cười quay sang, khẽ hỏi: ``\vie{Này, nếu em mỏi chân thì anh cõng em nhé?}''

Suzuki khựng lại một chút, lắc đầu ngay: ``\vie{Dạ không\dots\ em ổn ạ. Em\dots\ chưa cần đến mức đó đâu.}''

Cô bé chỉ cười nhẹ và bước tiếp. Tôi không gặng nữa. Có lẽ em ấy vẫn còn ngại, hoặc cũng có thể là vì điều gì khác.

Tôi chẳng thể biết hết ngay từ lần đầu, nhưng tôi chắc chắn một điều: trong cái thế giới đen đặc này, dù chỉ mới gặp nhau, tôi cũng không thể để cô ấy phải đi một mình, vì ai biết đâu được chuyện gì sẽ xảy ra bất ngờ ở một nơi như thế này?

\ct

\pov{Hikawa Suzuki}

Tôi đi bên cạnh anh ấy, một cách chậm rãi và từ tốn. Ban đầu là vì mỏi chân, nhưng càng đi tôi lại càng thấy mình muốn ở cạnh người này thêm một chút nữa.

Không phải vì tôi tin tưởng hoàn toàn, mà vì\dots\ hình như lâu lắm rồi tôi mới được nghe một giọng nói dịu dàng đến thế, không quát tháo, không cộc cằn, không ra lệnh, cũng không giễu cợt, như những học sinh nam khác mà tôi biết ở lớp.

Kuon-nii-san đi chậm lại, cùng bước với tôi. Tôi có thể thấy anh ấy liên tục để mắt tới tôi từ khóe mắt, như thể sợ tôi trượt chân hay mệt mỏi mà không nói ra. Nhưng anh không hỏi nhiều, cũng không tỏ ra vồn vã.

Mọi cử chỉ ấy\dots\ giống Onee-sama hay làm với tôi, hồi mà gia đình vẫn còn đông đủ.

Onee-sama của tôi là một người mạnh mẽ. Chị là người duy nhất từng đứng ra bênh vực tôi khi lũ bạn cùng lớp bắt nạt, là người nắm tay tôi chạy khỏi khu chung cư cháy hôm đó, là người luôn sẵn sàng đối đầu với bất kỳ ai dám tổn thương tôi.

Mặc dù\dots\ chị ấy có hay nóng giận, tất cả chỉ để bảo vệ tôi khỏi ``lũ dâm tặc dê xồm ngoài kia'' - theo lời chị ấy nói. Chị ấy dường như không tin tưởng một người đàn ông nào cả.

Tôi ngước nhìn Kuon-nii-san rồi cúi đầu xuống, hai tay đan vào nhau. Tôi cất tiếng, chậm rãi, vẫn còn chút run: ``\vie{Kuon-nii-san\dots\ Onee-sama\dots\ là người rất quan trọng với em.}''

Anh ấy ngoảnh đầu xuống dưới, nhìn tôi với đôi mắt có chút bối rối: ``\vie{Em rất quý trọng chị mình nhỉ.}''

Tôi gật đầu nhẹ, thốt lên một tiếng ``dạ'' rất bé. ``\vie{Em\dots không giỏi nói chuyện với người khác, nhưng chị ấy thì khác ạ. Chị ấy luôn bảo vệ em, luôn cho em ăn những món ngon, luôn để em chọn giường trước, và luôn cãi nhau với mẹ để bênh vực em mỗi khi em làm gì sai trái\dots}''

``\vie{Onee-sama của em hẳn phải là một người tuyệt vời lắm nhỉ,}'' anh ấy mỉm cười.

Tôi lại gật đầu lần nữa. Mắt tôi bắt đầu thấy cay cay, cổ họng nghèn nghẹn.

\vie{Chị ấy có thể hơi khó gần, hơi dữ, nhưng chị ấy không bao giờ đánh em hay mắng em vô cớ\dots\ Chị ấy yêu em lắm\dots\ em biết chứ.}

Tôi nhớ Onee-sama\dots\ Nhớ mà những lần chị ấy luôn bên cạnh, luôn than vãn mỗi khi đi làm về\dots\ Nhớ cả những lần mà hai chị em chơi chung với nhau mỗi khi chúng tôi rảnh\dots

Onee-sama\dots\ Chị đang ở đâu?

Tôi giương đôi mắt của mình lên, như muốn thỉnh cầu với anh ấy đầy tha thiết: ``\vie{Vậy nên\dots\ nếu anh có gặp chị ấy, nếu chị ấy có làm gì đó với anh\dots\ thì xin anh đừng trách chị ấy\dots\ Hãy coi đó như là một lời hứa với em đi.}''

Tôi cắn môi, ngẩng lên nhìn anh ấy lần nữa, mong rằng anh ấy có thể thấu hiểu.

Rồi, Kuon-nii-san mỉm cười rất nhẹ. ``\vie{Được rồi. Anh hứa.}

Nụ cười ấy không châm biếm, cũng không thương hại, mà giống như\dots\ một sự thông cảm. Tôi không giải thích được, nhưng tôi cảm nhận được như vậy.

Chúng tôi lại tiếp tục bước đi trong yên lặng, thỉnh thoảng hai chúng tôi có nói chuyện những thứ vẩn vơ để giết thời gian, từ lúc đó, tôi mới nhận ra rằng, dù là người xa lạ, Kuon-nii-san luôn giúp tôi ổn định lại tinh thần.

Dù thế, tôi luôn thắp lên một niềm hy vọng cho bản thân rằng tôi và Onee-sama có thể gặp lại nhau, dù rằng\dots\ tôi không biết liệu tôi có thể không.

Giờ đây, tôi biết mình không còn đơn độc nữa. Có một người đi ở bên cạnh, tôi cảm thấy an tâm phần nào.

Tôi quay sang anh ấy, mỉm cười dịu: \vie{Cảm ơn anh\dots\ vì đã đi cùng em, Kuon-nii-san\dots}''

\il

Tôi vẫn cố giữ khuôn mặt bình thường, dù có một nỗi khó chịu nhỏ đang âm ỉ trong người, không phải mệt mỏi, không phải lo sợ, mà là cảm giác cần nén lại một nhu cầu không tiện nói ra.

Tôi không để lộ điều đó. Đã quen rồi. Tôi từng phải giữ im lặng nhiều hơn thế khi còn ở chung cư. Chỉ là lần này, tôi không nghĩ nó sẽ kéo dài lâu như vậy.

Từng bước chân khiến tôi thấy khó chịu hơn mức bình thường. Một cảm giác âm ỉ mà tôi cố gắng quên đi suốt từ khoảng thời gian qua.

Tôi lén đưa tay về phía trước bụng dưới, rồi khẽ giữ chặt lại, kín đáo hết mức có thể. Tôi không muốn anh ấy biết. Tôi đã quen với việc này từ lâu rồi, quen với việc không nói ra, không than vãn, và không để lộ dù chỉ một chút. Ở khu chung cư cũ, từng giọt nước cũng là điều chân quý, nên đi vệ sinh thường xuyên bị xem là xa xỉ. Vậy nên tôi đã luyện được khả năng chịu đựng như thế.

Dẫu chẳng biết tôi có thể nhịn được bao lâu nữa. ``\vie{\hl{Ôi, mót quá\dots}}'' tôi khẽ lẩm bẩm, không để cho anh ấy có thể nghe thấy.

Kuon-nii-san chợt quay đầu lại, hỏi tôi có mệt không. Tôi thoáng giật mình, vội buông tay xuống và nở một nụ cười nhẹ: ``\vie{Dạ không sao đâu ạ, em vẫn ổn.}''

Chỉ là, tôi mong sớm tìm được Onee-sama. Và một nơi nào đó có thể gọi là ``\emph{an toàn}'', theo cả nghĩa đơn giản nhất.

\il

\pov{Hikawa Kuon}

Tôi im lặng suốt cả đoạn đường ngắn sau đó, chỉ lặng lẽ bước bên cạnh Suzuki-chan. Những lời mà cô bé vừa nói vẫn vang vọng trong đầu tôi, đặc biệt là ánh mắt khi em ấy nói ``đừng trách chị ấy''. Không hẳn là cầu xin, nhưng có một thứ gì đó giống như sự bảo vệ vô điều kiện dành cho người chị mà em tin tưởng nhất.

Giọng nói của em ấy nhỏ nhẹ, ngập ngừng, nhưng thật lòng. Nó không có gì màu mè khoa trương, không lời lẽ văn hoa, nhưng lại khiến tôi cảm thấy như mình đã bước vào một vùng đất mà từ lâu em ấy luôn giữ kín, không để ai tiếp cận tới.

Tôi nhớ lại người phụ nữ đã tấn công tôi trước đó, cú đấm mạnh đến mức làm bề mặt The Void nứt toác ra, ánh mắt sắc lạnh gần như điên dại, tôi mới bắt đầu suy ngẫm.

Chẳng lẽ đó là\dots\ chị của Suzuki-chan?

Cô gái mạnh mẽ ấy. Người mà cô bé gọi là ``Onee-sama.''

Một người con gái mà ngay cả khi em gái phải ở trong cảnh ngộ kỳ lạ thế này, vẫn để lại trong lòng em một tình thương đến mãnh liệt như vậy\dots\ thì dù trước đó có ra sao, tôi cũng không thể trách được.

Tôi quay sang nhìn Suzuki, cô bé vẫn cúi đầu lặng lẽ bước. Tôi mỉm cười dịu dàng, không lên tiếng một lời nào.

Có lẽ, từ giờ trở đi, tôi cũng sẽ học cách giữ lại một lời nào đó, như cách mà em vẫn đang làm.

Tôi hít một hơi thật sâu, rồi chậm rãi lên tiếng. ``\vie{Anh không biết mình sẽ đi đến đâu, và thú thực là anh cũng chưa từng bị vướng vào chuyện điên rồ nào như thế này\dots}''

Tôi nhìn về phía trước, nơi chỉ là một màu đen kéo dài vô tận. Tôi nhìn về Suzuki-chan, mỉm cười: ``\vie{Nhưng\dots\ nếu em muốn, anh sẽ đi cùng em tìm chị em. Anh không phải người xấu, và chắc chắn sẽ không làm gì tổn hại đến em. Anh hứa đấy.}''

Phải. Tôi đã từng trải qua những sự kiện tới mức phải gọi là gàn dở và kinh hoàng mà cả hai mươi lăm nồi bánh chưng tôi chưa thấy.

Một con bạch tuộc khổng lồ xuất hiện chỉ vì sự ham muốn với trò Spla***n\footnote{Xem lại {\flaregothic ÉPISODE 23}}.

Aku-tan bị mất trí nhớ ngay sát buổi biểu diễn solo đầu tiên của em ấy\footnote{Xem lại {\flaregothic SPÉCIAL \ab X}, quyển số Năm.}

Thậm chí là cả việc chúng tôi bị lạc vào một nơi gọi là ``Wonderland'' nhưng tràn đầy sự tạp nham\footnote{Xem lại {\flaregothic ÉPISODE 100}, quyển số Tám.}

Nhưng việc bản thâm một mình tôi bị kẹt vào một nơi đầy rẫy sự bí ẩn này thì tôi chưa từng bị như thế này bao giờ. Một danh sách ``những thứ kỳ quặc và mất não nhất mà Hikawa Kuon từng trải'' lại tiếp tục nới rộng.

Cô bé ngước nhìn tôi, đôi mắt ấy dường như đang mở to vì bất ngờ. Nhưng rồi, chỉ vài giây sau, ánh nhìn của em dịu lại. Em ấy nở một nục cười nhẹ, khuôn mặt nhẹ nhõm đi rõ rệt.

``\vie{\hl{\dots Em tin anh, Kuon-nii-san ạ.}}''

Tôi cảm thấy một điều gì đó ấm lên trong lồng ngực. Dù em chỉ nói mấy từ, nhẹ như gió, nhưng nó như thứ ánh sáng hiếm hoi trong cái thế giới tối tăm này.

``\vie{Vậy đi thôi,}'' tôi chìa tay ra. ``\vie{Cùng tìm người mà em quý nhất nào.}''

Suzuki chần chừ một chút, nhưng rồi cũng nhẹ nhàng đặt bàn tay nhỏ bé của em vào tay tôi. Bàn tay ấy lạnh, nhưng không còn run nữa.

Tôi mỉm cười. Không biết liệu mình có dẫn em đi đúng đường không, nhưng chí ít, từ giờ trở đi, chúng tôi sẽ không còn lạc lõng đơn độc trong cái cõi này nữa.

\end{document}